In a NoC router, multiple flits may request the same limited resource in the same cycle. For example, several input virtual channels may request the same output virtual channel, or several input ports may try to access the same output port of the switch. Since each resource can only be used by one flit or packet at a time, arbitration is required to decide which request should be served.

An arbiter is used when multiple requesters compete for one shared resource. It receives several request signals and grants the resource to one selected requester. In contrast, an allocator handles a more general case where multiple requesters compete for multiple resources. Therefore, an allocator can be viewed as a structure built from several arbiters to generate a conflict-free matching between requesters and resources.

In this router, allocation is mainly used in two places. The first one is virtual-channel allocation, where input virtual channels compete for available output virtual channels. The second one is switch allocation, where flits from different input sides compete for access to output ports through the switch. Since these allocation steps are on the critical path of router operation, the allocator should be simple, fast, and suitable for hardware implementation. Therefore, this design uses a separable input-first allocation method based on round-robin arbiters.
